\section{Audited routes and the remaining algebraic certificates}
\label{sec:root-routes}

\subsection{What a single fixed layout can never accomplish}

\begin{proposition}[RT-CROFTON: graph targets and fixed-layout limits]
\label{rt:crofton}
Let $P$ be a finite set of points in a Euclidean space, and let $C$ be the
product of cables of integer widths $c_{xy}\geq0$ between pairs in $P$.
Every finite bit network producing $C$ has Euclidean cable cost at least
$\sum_{xy}c_{xy}\|x-y\|_2$. The same conclusion holds if its output is a
coefficientwise limit of outputs with the same vertices and cable widths.
\end{proposition}
\begin{proof}
The proof of Proposition~\ref{prop:single-network} uses only that the
flattening rank of the target across each affine hyperplane is exactly
$2^{\sum_{xy\text{ crossing}}c_{xy}}$. This identity holds for every
product of independent wires. The network rank bound and its Crofton
integral therefore give the asserted inequality without a diameter or
unit-ball assumption. In a fixed-layout limiting family the flattening
rank is bounded by the same integer for every parameter. Every minor of
larger order is identically zero, and is therefore zero in the limit.
The same cut inequalities consequently hold at the limit. This argument
also works for a formal Laurent family with a well-defined constant limit:
the minors remain zero after reduction modulo the parameter.
\end{proof}

Thus an isolated border-rank or matrix-multiplication identity cannot
produce a Euclidean shortening when interpreted as one fixed cable
layout with the required graph target. A candidate sum of layouts must
be evaluated with its full selector cost. This does not rule out sums
with asymptotically negligible cover count, for which different terms
can have different deficient cuts.

\subsection{Finite price obstructions can be witnessed using wire targets}

\begin{proposition}[RT-FINITE: a complete-graph obstruction certificate]
\label{rt:finite}
Assume the all-tensor duality in Anthropic Theorem~A.6, pp.~20--21.
Let $w$ be a finite pseudometric, and identify its zero-distance points.
Write $W=\sum_e w(e)$ over unordered pairs of distinct quotient points.
The metric $w$ fails to be a price if and only if there are an integer
$m\geq1$, a number $\epsilon>0$, and covers of the complete graph tensor
$C_m^{\otimes N}$, for arbitrarily large $N$, with the following properties:
\begin{enumerate}
\item the number of covering tensors is $2^{o(N)}$;
\item every covering tensor is a product of binary wires on the same
parties, with $w$-weighted cable cost at most $mN(W-\epsilon)$.
\end{enumerate}
For the forward implication one may use products of $N$ members of one
fixed finite family, each member consisting of only a single cable.
\end{proposition}
\begin{proof}
The cases of fewer than two quotient points are trivial prices and have
no such certificate. Otherwise all $w(e)>0$. Let $K$ be the compact set
of prices $p$ with $p\leq w$ coordinatewise. It contains zero. If $w$
is not a price, compactness gives a number $\delta>0$ such that
$\sum_e p(e)\leq W-\delta$ on $K$: equality to $W$ under the coordinate
bounds would imply $p=w$.

Choose $m$ with $\max_e w(e)\leq m\delta/2$ and put
$q_e=\lceil m(W-\delta)/w(e)\rceil$. Let $A_e$ be the cable with
$q_e$ binary bits on pair $e$. For any spectral metric $d$, set
$M=\max_e d(e)/w(e)$. If $M>0$, then $d/M\in K$, so
$m\sum_e d(e)\leq m(W-\delta)M\leq\max_e q_e d(e)$.
If $M=0$, the same inequality is zero on both sides. Exponentiating
gives $\varphi(C_m)\leq\max_e\varphi(A_e)$ for every spectral point
on all tensors. Theorem~A.6 gives covers of $C_m^N$ by products of
$N$ members of this finite family, with cover count at most $2^{\eta N}$
eventually for every $\eta>0$. Choosing a sequence of such covers with
$\eta\to0$ gives the required subexponential count. Each factor costs
$q_e w(e)\leq m(W-\delta)+\max_e w(e)\leq m(W-\delta/2)$.
Use $\epsilon=\delta/2$.

Conversely, if $w$ is a price, Corollary~\ref{cor:price-normal-form}
represents it exactly as $w=s d_\psi$ for a normalized spectral point
and a finite $s$. Apply $\psi$ to the covers and take base-two logarithms:
$mNW\leq mN(W-\epsilon)+s\log_2 A_N$.
The last summand is $o(N)$, a contradiction.
\end{proof}

This is an operational reduction, not an explicit classification of
four- or five-party metrics. Its forward implication is nonconstructive:
it does not supply the covering coefficients needed for a concrete
shortening. It does specify precisely that the spectra used in the
separation are defined on all tensors; an extension from an assumed
wire-only functional is not being smuggled into the argument.

\subsection{Exact checks of the attached small border-rank curves}

\begin{proposition}[RT-BORDER: checked finite coefficient identities]
\label{rt:border}
The following explicit families in the supplied \texttt{border\_rank.tex}
have the stated Laurent coefficients:
\begin{enumerate}
\item In the subsection ``The nineteen pure terms explicitly'',
the seventeen-term determinant seed has no negative terms and constant
term $40\operatorname{Det}_3$ over
$\mathbb Z[z,h]/(z^4-z^2+1,h^2-5)$.
\item After both insertions in that subsection, the full nineteen-term
family has no negative terms and constant term forty times the ten
specified $2\times2$-minor slices. Dividing by forty requires
characteristic different from two and five, and the displayed
coefficients can be taken in an algebraic extension containing $z,h$.
\item The squarefree family in the subsection ``An eighteen-term curve
and three exact minima'' has eighteen pure terms and the ten specified
constant slices, over every characteristic.
\item The product constructions with respectively two and five graph
edges, or three and four graph edges, have eighteen and twenty pure
terms and respectively ten and twelve constant slices.
\end{enumerate}
These assertions concern the listed upper degenerations, not lower
bounds or the correctness of the entire attached document.
\end{proposition}
\begin{proof}[Exact certificate and mathematical verification]
The complete coefficient data are in
\path{inputs/rt_border_seeds.json}; the executable certificate is
\path{programs/rt_check_border_seeds.py}, and its recorded results
are in \path{outputs/rt_border_seeds.json}.
For the determinant seed the program represents each coefficient by
eight integers, for the basis $1,z,z^2,z^3,h,hz,hz^2,hz^3$, reducing
$z^4=z^2-1$ and $h^2=5$ exactly. It multiplies all ordered tensor
coordinates with nonpositive total weight and compares them with
$40\operatorname{Det}_3$. It then independently expands the actual
nineteen terms after transposition, cofactor signs, and the two added
minor insertions; every negative Laurent coefficient is checked to
vanish and every constant coordinate is checked against the ten
explicit minor matrices. There is no numerical root approximation.

For the eighteen-term squarefree curve, write $u,v$ for the squarefree
cell supports and $M$ for a target slice support. The coefficient is
$\sum_{u\cup v\subseteq P\subseteq M}(-1)^{|M|-|P|}$, equal to one
when $u\cup v=M$ and zero otherwise. The associated Laurent exponent
is $|u|+|v|-|M|$, which is $|u\cap v|\geq0$ when the coefficient is
nonzero. The program checks all $10\cdot11\cdot11=1210$ coordinates
and the eighteen-point down-set. For a graph edge $ij$ not meeting the
base index, the basic product-curve identity is
\[
 \frac{(e_0+t e_i+t e_j)^{\otimes2}-(e_0+t e_i)^{\otimes2}
 -(e_0+t e_j)^{\otimes2}+e_0^{\otimes2}}{t^2}
 =e_i\otimes e_j+e_j\otimes e_i.
\]
For an edge $0j$, use
$((e_0+t e_j)^{\otimes2}-e_0^{\otimes2})/t$,
whose constant term is the same symmetrized edge. Collecting repeated
point vectors and taking the indicated two products gives the stated
term counts and constant slices. All identities are over the integers,
so their reductions in characteristic two are also valid.
\end{proof}

None of these finite tensors has yet been embedded into a valid sum of
spatial wire layouts whose full cost beats the target. The missing
step is a target-preserving restriction with all auxiliary cable
widths and selector bits accounted for. A small border rank by itself
is not that step.

\subsection{The two-cut crossing problem and complex sums of squares}
\label{sec:root-sos}

For an integer $q\geq1$, let
\[
 T_q(a,b,c,d)=\langle a,c\rangle\langle b,d\rangle,
 \qquad a,b,c,d\in\mathbb C^q,
\]
where the pairing is bilinear, without complex conjugation. Define
$\chi_F(q)$ over a field $F$ to be the least possible value of
\[
 \max\{\rank_{ab\mid cd}S,\rank_{ad\mid bc}U\},
 \qquad T_q=S+U.
\]
Write $\sigma_{\mathbb C}(q)$ for the least number of squares of complex
bilinear forms whose sum is
$(\sum_i a_i^2)(\sum_j b_j^2)$. The connection with noncommutative
complexity is classical; see Hrube\v{s}, Wigderson, and Yehudayoff
(2011), Proposition~4.1, pp.~883--885. The following precise two-cut
form and finite geometric threshold are recorded because they are
the quantities used by the network search.

\begin{proposition}[RT-SOS: two-cut rank and a finite shortening threshold]
\label{rt:sos}
For every $q\geq1$,
\[
 \chi_{\mathbb C}(q)\leq\sigma_{\mathbb C}(q)
 \leq4\chi_{\mathbb C}(q).
\]
Over any field, absence of a planar strict shortening implies, for
every integer $m\geq1$,
\[
 \log_2\chi_F(2^m)\geq
 \max\{0,\sqrt2(m-1-\sqrt{2m-1})\}.
\]
Consequently an unbounded sequence of two-cut decompositions with
$\chi_F(q)\leq q^{\sqrt2-\epsilon}$, for any fixed $\epsilon>0$,
would give a planar strict shortening. Over $\mathbb C$ the same is
true with $\sigma_{\mathbb C}$ in place of $\chi_{\mathbb C}$.
\end{proposition}
\begin{proof}
Polarizing a sums-of-squares identity in each argument gives
\[
 2T_q(a,b,c,d)=\sum_{j=1}^{\sigma_{\mathbb C}(q)}
 \bigl(B_j(a,b)B_j(c,d)+B_j(a,d)B_j(c,b)\bigr).
\]
The two displayed sums have rank at most $\sigma_{\mathbb C}(q)$
across the two prescribed cuts. Conversely, write a decomposition
with both cut ranks at most $R$ as sums of $R$ separated products.
Setting $c=a,d=b$ expresses the product of the two quadratic forms
as a sum of at most $2R$ products of bilinear forms. Each product is
the sum of two complex squares, since
\[
 fg=\left(\frac{f+g}{2}\right)^2
      +\left(\frac{f-g}{2i}\right)^2.
\]
This proves the comparison, with no assertion that the same
polarization argument applies in characteristic two.

Now let $q=2^m$ and put $h=\log_2R$. For $h<\sqrt2m$, SH-12 gives
the exact minimum cost with real bridge width $h$ and four full
terminal cables of width $m$:
\[
 F(m,h)=\sqrt2\bigl(h+\sqrt{4m^2-h^2}\bigr).
\]
Suppose $F(m,h)+4<4m$. Expand $(S+U)^{\otimes t}$ into its $2^t$
term words. A word containing $p$ copies of $S$ can group them into
one adjacent-cut network with terminal widths $pm$ and bridge width
$\lceil ph\rceil$. Its other $t-p$ copies form a network for the
other adjacent cut, with terminal widths $(t-p)m$ and bridge width
$\lceil(t-p)h\rceil$. Take both symmetric optimal hub pairs, and
reorder the tensor factors locally at the terminals. Zero-width
groups are omitted. Thus every term uses at most four hubs and has
cost at most $tF(m,h)+4$: rounding costs less than two for each
of the two bridges, whose lengths are at most two. The selector
cost is exactly $4\log_2(2^t)=4t$. For
\[
 t>\frac{4}{4m-F(m,h)-4},
\]
this is an exact strict shortening with integer cable widths.

Therefore no planar shortening entails $F(m,h)+4\geq4m$.
For $m\leq3$ this gives no positive lower bound on $h$. For
$m\geq4$, solving this inequality on $0\leq h\leq\sqrt2m$
gives
$h\geq\sqrt2(m-1-\sqrt{2m-1})$; the second quadratic root
lies above $\sqrt2m$. If $h\geq\sqrt2m$, the same lower bound
holds automatically. This proves the finite estimate. For general
$q$, restrict all four index spaces to the preceding power of two.
Then $m=\lfloor\log_2q\rfloor$ and neither cut rank increases.
The finite bound is $\sqrt2m-O(\sqrt m)$, proving the last
assertions.
\end{proof}

\begin{proposition}[RT-SOS-UPPER: the checked crossing identity reaches the threshold]
\label{rt:sos-upper}
The local identities in Anthropic Lemma~B.1 and Proposition~B.2,
pp.~22--24, imply over every field that
\[
 \chi_F(q)\leq q^{\sqrt2+o(1)}\quad(q\longrightarrow\infty).
\]
If the Euclidean square is a price, then over $\mathbb C$,
$\sigma_{\mathbb C}(q)=q^{\sqrt2+o(1)}$.
\end{proposition}
\begin{proof}
Let $P(E)$ be the number of Frobenius symbols of energy at most
$E$. The cited counting lemma gives
$\ln P(E)=\kappa\sqrt E+o(\sqrt E)$, for a positive constant
$\kappa$. Among the at most $(E+1)^2$ classes of fixed size and
fixed energy at most $E$, let $Q(E)$ be the largest cardinality.
Then
\[
 \frac{P(E)}{(E+1)^2}\leq Q(E)\leq P(E),\qquad
 \ln Q(E)=\kappa\sqrt E+o(\sqrt E).
\]
Choose the least $E$ such that $Q(E)\geq q$, and a largest class
of energy $e\leq E$. Minimality and the asymptotic formula at
$E-1$ imply $\ln q=\kappa\sqrt E+o(\sqrt E)$.

Use this class on both crossing wires. Proposition~B.2 writes
$T_{Q(E)}=\sum_{N=0}^{4e}\Psi_N$. Group the terms with
$N\leq2e$ and those with $N>2e$. The first sum has rank at most
$(2e+1)^2P(2e)$ across $ad\mid bc$, by B.2(b); the second has
at most the same rank across $ab\mid cd$, by B.2(c).
Restrict the four index spaces to $q$ coordinates. The resulting
decomposition has maximum cut rank at most
\[
 (2E+1)^2P(2E)
 =\exp\bigl(\kappa\sqrt{2E}+o(\sqrt E)\bigr)
 =q^{\sqrt2+o(1)}.
\]
The decompositions are integral before reduction, so this holds
in every characteristic. If the square is a price, PG-SQUARE-GATE
and SH-15 exclude planar shortening. RT-SOS then gives the lower
exponent $\sqrt2$ over $\mathbb C$, while its factor-four
comparison supplies the matching upper bound for $\sigma$.
\end{proof}

The same conditional sums-of-squares exponent appears in Liu's
unrefereed manuscript (2026b, Section~2, ``Two crossing cuts,'' p.~3),
which cites Khoruzhii's crossing-decomposition upper bound (2026,
``Sum of squares'' section). These are related claims, not independent
verification of either VP--VNP manuscript. The explicit finite bound
in RT-SOS uses the exact cost function and tensor-word rounding above.
The established complex upper bound of Hrube\v{s} (2024,
Theorem~1, p.~12:2) has exponent $1.62$ and therefore does not
cross the required $\sqrt2$ threshold. None of these upper bounds
decides whether the square is a price.

There is an essential characteristic-two distinction: there
$\sigma_F(q)=1$, because
$(\sum_i a_i^2)(\sum_j b_j^2)=(\sum_{i,j}a_i b_j)^2$.
Hrube\v{s} (2024), Remark~2(i), p.~12:2, explicitly records this
degeneracy. It gives no crossing decomposition: the required
polarization divides by two. The invariant $\chi_F$ and the
integer-width geometric threshold remain meaningful in
characteristic two.

\subsection{What graph-image substitution does and does not certify}

\begin{proposition}[RT-IMAGE: exact residual images of graph tensors]
\label{rt:image}
Let $G$ be a finite graph with edge dimensions $q_e\geq1$, and let
$T_G$ be its product-of-wires tensor over $\mathbb C$. Partition its
vertices into $A$ and $B$, with at least two vertices in $B$, and let
$U$ be the image of the $A\mid B$ flattening, viewed as a tensor
space on the separate parties in $B$. Then
\[
 \dim U=\prod_{e\in\partial A}q_e,
 \qquad
 \min_{0\ne u\in U}\underline R(u)
       =\underline R(T_{G[B]}).
\]
Consequently the supplied image-substitution inequality gives
\[
 \underline R(T_G)\geq
       \prod_{e\in\partial A}q_e+
       \underline R(T_{G[B]})-1.
\]
For a fixed graph with binary edge loads $b_e$, recursively using
only these graph-image substitutions and elementary cut-rank base
cases cannot give an asymptotic border-rank exponent larger than
$\max_S\sum_{e\in\partial S}b_e$ for $T_G^{\otimes N}$.
\end{proposition}
\begin{proof}
Separate, at each vertex of $B$, its ports on the cut from its ports
internal to $B$. The flattening image is exactly
\[
 U=\left(\bigotimes_{v\in B}V_{\partial A,v}\right)
       \otimes T_{G[B]}.
\]
Indeed the internal wires on $A$ give a nonzero tensor on the
contracted side, and every cut port can be selected independently.
The cut-port tensor space has the displayed dimension. Every
nonzero element is $F\otimes T_{G[B]}$ for some nonzero boundary
tensor $F$. Choose one nonzero product-coordinate coefficient of
$F$, and contract the boundary ports with those coordinate
functionals locally at their respective vertices. The result is
a nonzero scalar multiple of $T_{G[B]}$. Restriction monotonicity
therefore gives $\underline R(F\otimes T_{G[B]})\geq
\underline R(T_{G[B]})$. A nonzero pure product $F$ attains equality.

Apply the image-substitution theorem in the supplied
\path{GCT_GINEQ.tex}, finding F-13, source lines~716--775. Its
proof was checked here: a $q$-dimensional limiting flattening image
meets the span of $r-q+1$ selected pure atoms in each rank-$r$
approximant, and a normalized limit gives the residual border-rank
bound. The graph residual calculation above is a distinct target
form of the coordinate-projection principle used in that source's
F-23, lines~886--943.

Each subsequent application removes a nonempty set of vertices and
adds one cut rank, subtracting one. There are at most $|V(G)|$
terms including an elementary terminal case. Every cut in an
induced subgraph has load at most the maximum cut load $M$ of
$G$: extend its two sides to a partition of all vertices; the
original cut edges remain and every additional load is nonnegative.
For $T_G^{\otimes N}$ every added term is thus at most $2^{NM}$.
The resulting numerical certificate is at most
$|V(G)|2^{NM}$, whose logarithm divided by $N$ tends to $M$.
This proves the scope limitation.
\end{proof}

The last assertion does not limit arbitrary substitution after
nongraph projections, independent stronger bounds for residual
tensors, or inequalities comparing two different graph tensors.
Those possibilities remain available. It identifies why repeated
application of the supplied image certificate to unchanged wire
graphs does not by itself decide the Euclidean-square gate.
\subsection{A characteristic boundary for the capacity route}
\label{sec:rt-capacity-characteristic}

\paragraph{RT-CAP-FIELD: the missing product step in characteristic two.}
The complex capacity obstruction is not transferred to $\mathbb F_2$
merely by observing that matrix multiplication is semistable there.
Christandl, Lysikov, and Zuiddam (2023), Lemma~23 and Example~24,
p.~312, give that semistability over arbitrary fields; their
Lemma~63, Appendix~A, p.~328, gives polystability by the same
irreducible-stabilizer criterion. The paragraph following their
Theorem~26, p.~312, distinguishes this fact from preservation of
semistability under Kronecker powers: it supplies the complex case
and leaves the arbitrary-field implication unproved.

The proof of PG-CAP-PRODUCT uses complex minimum-norm orbit
representatives for \emph{arbitrary sliced tensors}, not only the
matrix multiplication target. The cited arbitrary-field result does
not provide that product argument. Accordingly, the exact
left-versus-right flattening obstruction in SH-29 remains valid
over both requested fields, while PG-CAP-HYPER is asserted only
over $\mathbb C$. This records a failed transfer route; it is neither
a positive-characteristic counterexample nor a statement that no
different characteristic-two proof exists.

\subsection{An exact finite sums-of-squares check and its amplification limit}

\paragraph{RT-ZZ-CHECK.}
Zhang and Zhu (2026), equation~(3.2), p.~6, and Lemma~3.3,
p.~10 of arXiv version~2, give an explicit complex
$[12,12,18]$ composition formula. We independently transcribed
the displayed array and checked every polarized coefficient identity
\[
 v_{ij}\mathbin{\cdot}v_{k\ell}
 +v_{i\ell}\mathbin{\cdot}v_{kj}=2\delta_{ik}\delta_{j\ell}.
\]
The check uses the exact integer Gram matrix of their three exceptional
vectors and the fifteen orthonormal vectors, so no floating-point
arithmetic or conjugate inner product is involved.
All $12^4=20736$ identities pass. The input transcription is
\path{inputs/rt_zhang_zhu_12_12_18.json}; the independent program
\path{programs/rt_verify_zhang_zhu.py} writes
\path{outputs/rt_zhang_zhu_12_12_18_verification.json}.
This checks the positive formula only, not the cited real nonexistence
theorem or an asymptotic amplification.

Here is a precise obstruction to the naive amplification inference.
Tensoring bilinear composition maps and grouping their input spaces
does not in general produce another composition map. Already the
real two-square identity is complex multiplication on the basis
$1,i$. Its tensor square is multiplication in
$\mathbb C\otimes_{\mathbb R}\mathbb C$. In the orthonormal
coordinate basis $(1\otimes1,1\otimes i,i\otimes1,i\otimes i)$,
put
\[
 x=1\otimes1+i\otimes i,\qquad
 y=1\otimes1-i\otimes i.
\]
Their coordinate squared norms are both two, but
$xy=1\otimes1-(i\otimes i)^2=0$. Thus the proposed composition
identity would read $4=0$. The same integer coefficient failure
persists after extension to $\mathbb C$, using the bilinear
quadratic form rather than a Hermitian norm. Consequently the
finite $[12,12,18]$ formula alone gives no bound
$\sigma(12^N)\le18^N$. A valid additional composition or cover
construction is still needed for the asymptotic gate.

\paragraph{RT-BALANCED-SEARCH: normalized dense searches.}
We revisited the binary balanced $K_{3,3}$ conversion with one and
two summands, using dense complex local maps. In contrast to the
earlier exploratory parameterization, every local map and the vector
of summand coefficients is normalized inside the objective, with the
normalization derivative included. This removes a possible stopping
artifact caused by arbitrarily large raw parameter norms.
The grouped source and target are reconstructed by the graph routine
in \path{programs/sh_border_k33_search.py}; the optimization and its
independent finite-difference gradient check are in
\path{programs/rt_balanced_cover_search.py}.
Twelve starts per term count, with limits of $2000$ and $2500$
iterations, reached best normalized squared overlaps $1/4$ and $1/2$,
respectively. Gradient discrepancies were below $4\cdot10^{-12}$.
Inputs and complete outputs, including the best dense maps, are
\path{inputs/rt_balanced_q2_a1.json},
\path{inputs/rt_balanced_q2_a2.json}, and the matching
\path{outputs/rt_balanced_q2_a1} and
\path{outputs/rt_balanced_q2_a2} files.
No exact identity or degeneration was found. These local-search
values are not certified global optima and do not prove an obstruction
to any number of terms or any wire dimension.

\subsection{A smaller width-six search and the role of odd internal degrees}

\begin{lemma}[RT-CUT-SUBMOD: a stronger planar capacity bound]
\label{rt:cut-submod}
Place four terminals $a,b,c,d$ cyclically at the coordinate
crosspolytope in $\mathbb R^2$. For a geometric cable network let
$c_i$ be its singleton minimum-cut capacities, and let $H\ge L$
be its two adjacent-pair minimum-cut capacities. Its cost satisfies
\[
 \cost\ge(\sqrt2-1)\sum_i c_i+(2-\sqrt2)(H+L)
 \ge\sqrt2 H+(2-\sqrt2)L.                 \tag{RT-W6.1}
\]
The assertion allows arbitrary nonnegative real capacities.
\end{lemma}
\begin{proof}
The planar Crofton argument in (S3) applies directly to minimum
cut capacities: every separating line cuts edges of total capacity
at least the corresponding terminal minimum cut. This gives the
first inequality with the same singleton and adjacent-pair weights.
The union of cuts isolating two terminals separates their pair;
therefore each adjacent-pair minimum cut is at most the sum of
the two corresponding singleton minimum cuts. Applying this to
both sides of the cut of capacity $H$ yields $\sum_i c_i\ge2H$.
Since $2(\sqrt2-1)+(2-\sqrt2)=\sqrt2$, the second inequality follows.
\end{proof}

\begin{proposition}[RT-W6-PARAM: remaining parameters at width six]
\label{rt:width-six-parameters}
Over any field, a strict shortening of width $m=6$ must have
either $A=2$ and $2\le D\le19$, or $A=3$ and $2\le D\le16$.
\end{proposition}
\begin{proof}
The one-term case is excluded by Crofton. The continuous
necessary condition from SH-23 is
\[
 \frac6{\log_2A}>\gamma_D,
 \qquad\gamma_D=2+\sqrt2-\frac{1+\sqrt2}{D},
                                                        \tag{RT-W6.2}
\]
and $\gamma_D$ increases with $D$. For $A\ge7$, use
$\log_27>11/4$ (since $7^4>2^{11}$) and
$24/11<(3+\sqrt2)/2=\gamma_2$; the last inequality follows
from $225<242$. Thus $A\le6$.

For $A=5$, $\log_25>23/10$ because $5^{10}>2^{23}$.
Also $60/23<\gamma_3=(5+2\sqrt2)/3$, by $4225<4232$.
Hence only $D=2$ remains. Its cost is less than
$24-4\log_25<74/5$. Each adjacent cut has integer
capacity at most $10$, by the $\sqrt2$ path-length bound,
and the sum of the two rank ceilings is at most $14$.
Consequently one term contributes at most
$2^{10}+2^4=1040$ to the sum of its two adjacent ranks.
Five terms give at most $5200$, below the target sum
$2\cdot2^{12}=8192$.

For $A=6$, $\log_26>5/2$, so (RT-W6.2) again leaves
only $D=2$. Its cost is less than $14$, giving adjacent
capacities at most $9$; six ranks of at most $2^9$
cannot sum to the target rank $2^{12}$.

For $A=4$, $\gamma_6>3$ follows from $50>49$, so
$D\le5$. In dimensions $D=3,4,5$ the strict cost bound
is $w<8D$. The path-length bound gives capacities at most
$K=16,22,28$, respectively, and in each case $K=6D-2$.
The target rank on every sign cut is $2^{6D}=4\cdot2^K$.
All four term ranks must therefore equal $2^K$ on every
sign cut. The sign-frame bound then gives cost at least
$2K>8D$ for each term, a contradiction.

It remains to exclude $A=4,D=2$. Here $w<16$,
each adjacent rank ceiling is at most $11$, and their
sum in each term is at most $15$. A term whose larger
ceiling is at most $10$ has adjacent-rank sum at most
$2^{10}+2^5=1056$; every other term has sum at most
$2^{11}+2^4=2064$. Since $1056+3\cdot2064<8192$,
every term has larger ceiling exactly $11$. Its high
rank is greater than $2^{10}$. Each terminal space
has dimension $64$, and a pair flattening has rank
at most $64$ times either constituent singleton rank.
Thus all four singleton rank ceilings are at least $5$.
If the lower adjacent ceiling were at least $3$,
integer-ceiling Crofton would give
\[
 20(\sqrt2-1)+14(2-\sqrt2)=8+6\sqrt2>16.
\]
Hence every term has lower adjacent rank at most $4$.
Fix the two endpoints of one diameter to matching basis
vectors. The remaining target is the $64\times64$
identity matrix, while the slice of each summand has
rank at most either of its adjacent flattening ranks.
It follows that $64\le\sum_{j=1}^4\min(r_{j,1},r_{j,2})
\le16$, the final contradiction.

Finally the complement-even variance and Cantelli
condition from SH-26 requires
\[
 (D-1)^2(\log_2A)^2<(A-1)18D.
\]
For $A=2$ it fails at every $D\ge20$.
For $A=3$, use $\log_23>19/12$ to see that it fails
at every $D\ge17$, since
$361(D-1)^2\ge5184D$ in that range.
\end{proof}

\begin{proposition}[RT-W6-EULER: odd internal degrees are necessary]
\label{rt:width-six-euler}
If a strict shortening with $m=6,D=2,A=2$ exists,
\emph{each} of its two term networks has an internal vertex
whose total incident integer cable width is odd.
\end{proposition}
\begin{proof}
Write the two tensor summands as $S,U$. Their adjacent
rank ceilings are at most $12$, and the two ceilings
in each term sum to at most $19$, because $w<20$.
If on one adjacent cut both ceilings were at most $11$,
both ranks would have to equal $2^{11}$ to reach the
target rank $4096$. On the other cut some ceiling is
at least $11$, contradicting that term's sum bound.
Thus each cut has a ceiling $12$, and these belong to
opposite terms. Each low ceiling is at most $7$.
Each high rank is consequently at least $4096-128=3968$,
so every singleton rank ceiling of both terms is $6$.

Suppose one term network is \emph{inner Eulerian},
meaning that every internal vertex has even total
incident integer capacity. The four singleton terminal
cuts and the two adjacent-pair cuts have no pairwise
crossing triple. The integral multiflow locking theorem
therefore applies: Hirai (2010), Theorem~1.1, author
preprint p.~2, restating the Karzanov--Lomonosov theorem.
It gives an integral collection of terminal-to-terminal
paths simultaneously attaining all these minimum cuts.
Straighten every path to its terminal segment. Capacity
feasibility and the triangle inequality show that the
resulting integral weighted terminal graph has cost no
larger than the original network, and preserves the six
locked terminal cut values.

Its four terminal degrees are at least $6$ and its high
adjacent cut $H$ is at least $12$. Suppose its low cut
$L$ were at least $5$. Let $y$ be the total capacity of
the two diameter edges. The total adjacent-edge capacity
is $H+L-2y$, so its exact cost is
\[
 c=\sqrt2(H+L)+(2-2\sqrt2)y,
 \qquad y\le L.
                                                        \tag{RT-W6.3}
\]
If $H\ge13$, this is at least $10+8\sqrt2>20$;
if $L\ge6$, it is at least $12+6\sqrt2>20$.
Hence $H=12,L=5$. If $y\le4$, (RT-W6.3) is at least
$8+9\sqrt2>20$, leaving only $y=5$. The two adjacent
edges within the low-cut sides then vanish, and the
other two adjacent capacities sum to $7$. The total
terminal degree is $2(H+L-y)=24$, so all four degrees
are exactly $6$. Opposite degree equations force those
two adjacent capacities to be equal, hence both $7/2$,
contradicting integrality. We conclude $L\le4$.

Thus the low adjacent rank of the inner-Eulerian term
is at most $16$. The mixed-cut slice inequality for
the two summands forces the other term's low rank to
be at least $64-16=48$, with ceiling at least $6$.
Its high rank is at least $4096-16=4080$, giving high
ceiling $12$ and all singleton ceilings $6$. Crofton
now gives that term cost at least
\[
 24(\sqrt2-1)+18(2-\sqrt2)=12+6\sqrt2>20,
\]
the required contradiction. No parity restriction on
general tensor-network vertices is assumed; the result
identifies a necessary feature of any remaining candidate.
\end{proof}

The exact exploratory program
\path{programs/rt_planar_rank_budget.py} applies (RT-W6.1)
to integer rank-ceiling profiles. It uses outward rational
bounds on $\sqrt2$ and $\log_2A$, with input
\path{inputs/rt_planar_width6_rank_budget.json} and output
\path{outputs/rt_planar_width6_rank_budget.json}. It independently
excludes its $A=4,5,6$ planar relaxations, but admits four
$A=3$ profile multisets and one $A=2$ multiset. The latter
are feasible necessary data, not tensor constructions; the
analytic arguments above carry the stated universal conclusions.

These last feasible relaxation profiles are historical: the later
strict crossing-rank and parity arguments in
Corollary~\ref{cor:planar-width-six-seven} exclude both the two-term
and the three-term planar width-six cases. No tensor realizing
the earlier necessary data was constructed.

An additional numerical search allowed eight, ten, or twelve
internal vertices, six initializations per size, and five
alternations between integer capacities and vertex positions.
The input is \path{inputs/rt_integer_network_large.json}; the
program is \path{programs/sh_integer_network_search.py}; the full
output and trace are \path{outputs/rt_integer_network_large.json}
and \path{outputs/rt_integer_network_large_trace.jsonl}.
Its 47 accepted iterates had best cost approximately
$20.348984679394846$. This finite search supplies no lower bound;
the exact parity calibration in
Theorem~\ref{thm:critical-parity-calibration} supplies the
universal geometric obstruction.



\subsection{An integral rank witness and the scope of classical coding}

\begin{proposition}[RT-UNIVERSAL-RANK-WITNESS]
\label{rt:universal-rank-witness}
The eight-point rank-distance metric used in SH-20 has a
realization in $\operatorname{Mat}_3(F)$ over every field $F$.
The same integral matrices work in every characteristic.
Consequently its square-root metric is not conditionally
negative definite over any of these realizations.
\end{proposition}
\begin{proof}
Put $A=E_{11}$, $B=E_{22}$ and
\[
 C=\begin{pmatrix}1&0&-1\\-1&0&1\\1&1&0\end{pmatrix},
 \quad R=\begin{pmatrix}0&0&0\\1&0&-2\\0&0&0\end{pmatrix},
 \quad S=\begin{pmatrix}0&1&2\\0&0&0\\0&0&0\end{pmatrix}.
\]
For binary labels $x,y,z\in\{0,1\}$ define
\[
 P_{xy0}=xA+yB,\qquad P_{xy1}=C+xR+yS.
\]
The labels are ordinary integers before their images are taken
in $F$; no embedding of $\mathbb F_2$ into $F$ is being asserted.
Within either layer, the rank of a difference is the Hamming
distance of the two binary labels. In the first layer this
follows from the diagonal form. In the second, a nonzero
coefficient of $R$ or $S$ gives a unit entry, and when both
coefficients are nonzero the first two rows and columns have
unit determinant.

For a cross-layer pair the difference is
\[
 P_{xy1}-P_{uv0}=
 \begin{pmatrix}
 1-u&y&-1+2y\\
 -1+x&-v&1-2x\\
 1&1&0
 \end{pmatrix}.
\]
Its determinant, as an integer polynomial, is
\[
 (u+x-2ux)-(v+y-2vy).
\]
For binary arguments the two parenthesized expressions are
exactly the indicators $u\ne x$ and $v\ne y$. Thus the
determinant is zero if these indicators agree, and is $1$ or
$-1$ otherwise. The minor in rows $1,3$ and columns $1,3$ is
always $1-2y\in\{1,-1\}$. Hence every cross-layer rank is
exactly two or three, respectively, over every field.

In the label order $x+2y+4z$, the rank of a difference therefore
has the same translation-invariant table as SH-20:
\[
 (0,1,1,2,2,3,3,2).
\]
Its parity-character quadratic form is consequently
$8(2\sqrt3-\sqrt2-2)>0$, as computed in SH-20. This proves
the last assertion without changing the underlying metric.
\end{proof}

The initial finite search in
\path{programs/rt_rank_witness_odd_field.py} found a realization
over $\mathbb F_3$; its complete diagnostic is
\path{outputs/rt_rank_witness_F3.json}. Replacing the residues
$2$ in that realization by $-1$ led to the integral formula
above. The universal assertion is verified independently of
that search by the determinant calculation in the proof and by
\path{programs/rt_verify_universal_rank_witness.py}, with input
\path{inputs/rt_universal_rank_witness.json}. Its output
\path{outputs/rt_universal_rank_witness.json} supplies a unit
minor and the vanishing larger minors for every one of the
28 unordered pairs. These integer certificates include every
positive characteristic, with no exceptional primes.

\paragraph{Consequence for the causal-linear-code route.}
Apply the rank-metric inequality of EX-38 over the coding
field, using the matrices in Proposition~\ref{rt:universal-rank-witness}.
Its concave-rank version shows that every causal linear-code
wire inequality is satisfied by this same non-CND square-root
metric. This statement concerns the coding field, which need
not equal the field of coefficients of a tensor restriction.
Thus changing a linear code from a binary alphabet to an
alphabet over $\mathbb F_9$, or to another finite field, cannot
make its direct wire inequality exclude this eight-point
witness. This closes this particular route to the SH-22/SH-33
contraction gate; it does not apply to arbitrary linear maps
in tensor networks or to sums of restricted tensors.

\paragraph{Recent network-coding claims and the separate Euclidean barrier.}
Braverman and He (2026, Theorem~1.1, p.~1; deterministic
construction, \S4, pp.~9--11) report an undirected coding gain
over fractional routing. Their field alphabet is $\mathbb F_9$.
The full code and its finite data are not adopted here as a
checked premise. The distinction between a routing-cost
witness and a Euclidean witness is already decisive for its
direct use: any one-term graph-tensor restriction satisfies
every cut-rank inequality. Taking nonnegative combinations of
these inequalities preserves every finite $L_1$ wire metric,
and in particular Euclidean distance by the Crofton argument
of SH-01. A coding gain against some shortest-path metric
therefore supplies no Euclidean shortening by itself. The
stronger universal rank-witness argument above also rules out
a direct application of causal linear codes to the particular
non-CND configuration used by the present contraction gate.


\begin{proposition}[RT-STABILIZER-SELECTOR: a quadratic selector barrier]
\label{rt:stabilizer-selector}
Consider complex qubit stabilizer states, with local operations
consisting of party-local Clifford gates, local stabilizer
ancillas, and nonzero postselected Pauli measurements.
Let $4\mid r$ and $N\ge1$. Any conversion
\[
 G_r(2)^{\otimes N}\otimes\mathrm{GHZ}_{2}^{\otimes b}
 \longrightarrow H_r(2)^{\otimes N}
\]
using these operations requires
\[
                         b\ge \frac{Nr^2}{6}.
 \tag{RT-STAB-selector}
\]
The GHZ factors are shared among all $2r$ parties.
In particular, this stabilizer construction cannot supply the
subquadratic selector exponent required by SH-33.
This is a restriction on the specified realization of the local
maps, not on all complex tensor restrictions or all tensor covers.
\end{proposition}
\begin{proof}
Let $\mathcal S$ be the binary stabilizer space of a pure
stabilizer state, of dimension $n$. For each party $v$, the
local symplectic pairing restricted to $\mathcal S$ has matrix
$C_v$ in a common basis. These are the commutation matrices
of Englbrecht, Kraft, and Kraus (2022, Definition~5 and
Eqs.~(22)--(23), printed p.~8). They are alternating and
$\sum_v C_v=0$. Their behavior under party-local Clifford
changes is described by the same source, Theorem~7,
printed p.~10. We use the following rank consequence, including
its behavior under measurements, which is proved here.

Fix arbitrary $d$-by-$d$ matrices $P_v$ over $\mathbb F_2$,
and put
\[
 \mathcal R_P(\mathcal S)
   =\operatorname{rank}_{\mathbb F_2}
        \left(\sum_v P_v\otimes C_v\right).
 \tag{RT-STAB-rank}
\]
A common change of stabilizer basis acts by invertible left and
right factors on this matrix, so its rank is well defined.
Party-local Clifford gates preserve all local symplectic
pairings and hence preserve this rank. Adding a local pure
stabilizer ancilla adds a common radical summand to every
$C_v$, so it leaves the rank unchanged.

For a nontrivial Pauli measurement supported at a single party,
let $m$ be its binary Pauli vector. The old stabilizer first
restricts to the codimension-one subspace
$\mathcal S_0=\{s\in\mathcal S:\omega(m,s)=0\}$ and then adds
$m$. Because $m$ is supported at only one party, it has zero
local pairing with every vector of $\mathcal S_0$ at every
party. Thus the new tuple is a simultaneous restriction of the
old tuple followed by adjoining one common radical vector.
Rank cannot increase under this operation. If the measurement
was already fixed by the stabilizer, a nonzero outcome changes
nothing. The signs of the chosen nonzero outcomes do not change
the binary spaces. Discarded local systems may be measured and
conditioned further without changing a pure target, so the
same monotonicity applies to such implementations.

The rank in (RT-STAB-rank) is additive on tensor products,
because the stabilizer spaces and all their local forms become
direct sums. For one EPR pair on $uv$, use generators $X_uX_v$
and $Z_uZ_v$. Both endpoint forms are
$J=\left(\begin{smallmatrix}0&1\\1&0\end{smallmatrix}\right)$.
Consequently an EPR graph with multiplicities $a_{uv}$ has
\[
 \mathcal R_P=2\sum_{uv}a_{uv}
                      \operatorname{rank}(P_u-P_v).
 \tag{RT-STAB-wire}
\]
For a GHZ state on $M$ parties, use $X_1\cdots X_M$ and
$Z_1Z_i$ for $i=2,\ldots,M$. After block reordering, its
weighted matrix has the form
\[
 \begin{pmatrix}
 0&P_2-P_1&\cdots&P_M-P_1\\
 P_2-P_1&0&\cdots&0\\
 \vdots&\vdots&&\vdots\\
 P_M-P_1&0&\cdots&0
 \end{pmatrix}.
\]
The top block row has rank at most $d$ and the remaining block
rows together have rank at most $d$, so its rank is at most
$2d$. Combining this bound with (RT-STAB-wire) and monotonicity
gives the necessary inequality
\[
 \sum_{uv\in H}a_{uv}\operatorname{rank}(P_u-P_v)
 -\sum_{uv\in G}a_{uv}\operatorname{rank}(P_u-P_v)
 \le db.
 \tag{RT-STAB-cost}
\]

Now set $d=3$ and use the eight matrices of
Proposition~\ref{rt:universal-rank-witness}, reduced modulo two.
Place the four even-parity labels $0,3,5,6$ on the left
and the four odd-parity labels $1,2,4,7$ on the right.
The sum of rank distances over the sixteen cross pairs is
$24$: each even label has two rank-one and two rank-two
neighbors. Within either parity class, each label has one
rank-two and two rank-three neighbors, so the sum over the
six unordered pairs is $16$. The two target cliques therefore
have total rank cost $32$, with defect $32-24=8$.

For $r=4t$, repeat each label $t$ times. Distances between
identical labels vanish, and every other pair contribution is
multiplied by $t^2$. With $N$ wire copies, the defect is
$8Nt^2=Nr^2/2$. Substitution into (RT-STAB-cost) yields
$Nr^2/2\le3b$, proving the proposition.
\end{proof}

The same argument permits $o(N)$ additional source copies:
for fixed $r$, their rank-cost contribution is $o(N)$, so they
do not remove the positive leading lower bound on $b/N$.
It is essential that the entire GHZ-assisted conversion have
a stabilizer implementation. A sum of individually chosen
stabilizer local maps may require non-Clifford controlled maps
when the summands are selected coherently; the proposition
does not bound those general covers.


\subsection{A local-unitary counterexample to the weighted stabilizer rank}
\label{rt:lu-rank-counterexample}

\begin{proposition}\label{rt:lu-rank-failure}
The weighted binary rank $\mathcal R_P$ in (RT-STAB-rank) is not
invariant, and is not a monotone, under arbitrary party-local unitaries
between complex pure stabilizer states. There are two locally unitary
27-qubit stabilizer states and explicit matrices $P_v\in
\operatorname{Mat}_6(\mathbb F_2)$ of rank at most one for which the
ranks are respectively $36$ and $37$.
\end{proposition}
\begin{proof}
The state data are from Ji, Chen, Wei, and Ying (2010, preprint
pp.~9--10, Eqs.~(18)--(21)); only their displayed data, and not their
local-Clifford inequivalence theorem, are used. Let $G$ have the six
binary rows
\[
\begin{array}{l}
100010001010101010100011110\\
101010111001100000001010101\\
011001100111100111100110011\\
000111100000011001100001111\\
000000011111111000011111111\\
000000000000000111111111111.
\end{array}
\]
Let $S$ be its row space, and define $q(x)=\sum_{ij\in E}x_ix_j$
over $\mathbb F_2$, with the one-based edge set
\[
\begin{split}
 E=\{&(1,2),(1,3),(1,8),(2,4),(2,8),(2,16),\\
     &(3,4),(3,8),(3,16),(4,8),(8,16)\}.
\end{split}
\]
Write $J$ for this graph's symmetric binary adjacency matrix. The
27 local phase exponents are
\[
\begin{split}
(e_1,\ldots,e_{27})=(&3,5,7,5,1,3,5,7,1,3,5,5,3,7,\\
                  &3,3,7,1,7,3,1,5,5,5,3,5,3).
\end{split}
\]
For each of the $64$ elements $x\in S$, direct integer evaluation gives
$\sum_v e_vx_v\equiv4q(x)\pmod8$. Consequently
\[
 \bigotimes_{v=1}^{27}\operatorname{diag}(1,e^{\pi i e_v/4})
 \left(2^{-3}\sum_{x\in S}|x\rangle\right)
 =2^{-3}\sum_{x\in S}(-1)^{q(x)}|x\rangle.
 \tag{RT-LU-phase}
\]
This is an exact invertible party-local unitary conversion.

Here is complete finite data for the new weighted-rank witness. Put
\[
 A=\begin{pmatrix}
1&1&0&0&1&1\\1&1&1&1&0&1\\0&1&1&1&0&0\\
1&1&0&1&1&1\\1&0&1&0&0&1\\1&1&1&1&1&1
\end{pmatrix},\qquad
 B=\begin{pmatrix}
1&0&0&1&0&0\\0&1&1&0&1&0\\1&1&1&1&0&1\\
1&0&0&1&1&1\\1&1&1&0&0&0\\1&0&0&0&0&1
\end{pmatrix}.
\]
If $g_v$ is column $v$ of $G$, take
$P_v=(Ag_v)(Bg_v)^{\mathsf T}$. Three of these matrices are zero;
all have rank at most one. Choose any $21$-row basis $H$ of
$S^\perp$. Binary generator matrices for the two stabilizer spaces are
\[
 X=\begin{pmatrix}G\\0\end{pmatrix},\quad
 Z^{(0)}=\begin{pmatrix}0\\H\end{pmatrix},\quad
 Z^{(1)}=\begin{pmatrix}GJ\\H\end{pmatrix}.
\]
Each pair $(X,Z^{(i)})$ has independent rows and zero global
symplectic pairing: $GH^{\mathsf T}=0$ and $GJG^{\mathsf T}$ is
symmetric. The corresponding signed generators stabilize the two
states in (RT-LU-phase). Generator signs do not affect the local
commutation forms
\[
 C_v^{(i)}=X_{:v}(Z^{(i)}_{:v})^{\mathsf T}
             +Z^{(i)}_{:v}X_{:v}^{\mathsf T}.
\]
Exact elimination over $\mathbb F_2$ on the two explicitly specified
$162$-by-$162$ matrices gives
\[
 \operatorname{rank}\sum_{v=1}^{27}P_v\otimes C_v^{(0)}=36,
 \qquad
 \operatorname{rank}\sum_{v=1}^{27}P_v\otimes C_v^{(1)}=37.
 \tag{RT-LU-ranks}
\]
This finite calculation is included as part of the proof:
\texttt{rt\_lu\_rank\_probe.py} supplies the initial exact computation;
\texttt{pg\_verify\_lu\_rank\_counterexample.py} independently constructs
$S^\perp$ with different pivots, builds the local forms entry by entry,
and uses plain-list binary elimination. Both also check all $64$
phase congruences. The compact matrices above are recorded in
\texttt{inputs/rt\_lu\_rank\_witness.json}; their expanded $P_v$ agree
entry by entry with the independently verified data in
\texttt{outputs/rt\_lu\_rank\_probe.json}. Thus the local-unitary
conversion increases $\mathcal R_P$, proving the assertion.
\end{proof}

The proposition closes an attempted extension of
Proposition~\ref{rt:stabilizer-selector}: dropping the restriction to
local Clifford gates and Pauli postselection would invalidate its
rank-monotonicity step. It does not invalidate the proved restricted
barrier, and gives neither a general graph-tensor conversion nor a
strict shortening.


\subsection{Further operator identities: a precise remaining lift}
\label{rt:operator-cauchy-scope}

Itoh (2012, author preprint pp.~24--27, Proposition~6.6,
Eq.~(6.3), and Theorem~6.10) retains a group-algebra-valued
preimmanant rather than applying a character immediately.
However, the Cauchy expansion in Eq.~(6.3)/Theorem~6.10 takes
values in the \emph{center} $Z\mathbb C S_p$:
\[
 \operatorname{preimm}_p X
 =\sum_{\lambda\vdash p}s_\lambda(a_1,\ldots,a_n)
                       \widetilde s_\lambda,
 \qquad
 \widetilde s_\lambda=\frac1{p!}
       \sum_{\sigma\in S_p}\chi_\lambda(\sigma)\sigma.
\]
These central elements specify scalar actions on irreducible
symmetric-group factors. They do not supply the noncentral matrix
entries between local multiplicity vectors that a full tensor lift
would require. The cited Cauchy--Binet formula instead concerns
matrix multiplication with commuting entries; no correspondence
from its indices to party-local maps of the matching-deleted graph
has been established here. Thus this source is not being used as
a proof of the required cover.

A more general lift using several dense local-operator products
per matching remains open. One concrete version is to split each
local $(\mathbb C^2)^{\otimes(r-1)}$ into its Schur--Weyl sectors
and seek simultaneous maps on the multiplicity spaces. There are
only polynomially many sector labels at one party, but the
multiplicity spaces themselves can have exponential dimension.
Counting sector labels alone therefore does not bound the number
of product maps needed to reconstruct a sector. The all-$r$
one-product-per-matching obstruction in
Proposition~\ref{pg:cauchy-lift-nogo} applies to the stated
quotient-preserving maps, and does not dispose of this multiple-term
problem. No improvement to the balanced-cover threshold is inferred.

\subsection{A quantitative remaining gap in the recoupling route}

\begin{proposition}[RT-OP-GAP: failure forces an exponentially small gap]
\label{rt:op-gap}
Use the real orthogonal realization $V=[\lambda]$ of
Proposition~\ref{pg:op-sos-duality}, and write $M=\dim V$.
Let $E$ be the real symmetric commutant on $V\otimes V$, with
Hilbert--Schmidt inner product, and let $\Gamma$ be partial transpose.
Let $L\subset E$ consist of the operators supported on the isotypic
summands with entropy below $h$, and let $\Pi:E\to L$ be orthogonal
projection. Put
\[
 B=\sum_{H(\nu/n)<h}d_\nu g_{\lambda\lambda\nu},\qquad
 A=\Pi\Gamma\Pi\big|_L,\qquad T=I_L+A.
 \tag{RT-OP-gap-data}
\]
Suppose that no $X\in L^\perp$ satisfies $\Gamma X=X$ and
$\operatorname{Tr}X\ne0$. Then $B>0$, $T$ has a positive eigenvalue,
and its smallest positive eigenvalue $\delta$ obeys
\[
 \delta\leq\frac{2B}{M^2}.
 \tag{RT-OP-gap}
\]
For a maximum-dimension or rectangular $q$-row type as above, this
implies
\[
 \delta\leq
 2(n+1)^{B_q+2A_q}\,2^{-n(2\log_2q-h)}.
 \tag{RT-OP-gap-asymptotic}
\]
Thus a lower bound on $\delta$ larger than this quantity would supply
the required high-support operator. No such uniform lower bound is
proved here.
\end{proposition}
\begin{proof}
All orthogonal complements are taken within $E$. Partial transpose
is an orthogonal involution of $E$ and fixes the identity $I$ on
$V\otimes V$. Write $E_+$ and $E_-$ for its two eigenspaces.
The assumption is that $I$ is orthogonal to $L^\perp\cap E_+$.
Finite-dimensional orthogonal complementation gives
\[
 I\in (L^\perp\cap E_+)^\perp=L+E_-.
\]
Consequently there is $Y\in L$ with
\[
 Y+\Gamma Y=2I.
 \tag{RT-OP-gap-equation}
\]
Conversely this equation would make every such high-support
$\Gamma$-fixed $X$ traceless, by taking its inner product with $X$.

For $Z\in L$, orthogonality and involutivity of $\Gamma$ give
\[
 \langle Z,TZ\rangle
 =\langle Z,(I+\Gamma)Z\rangle
 =\tfrac12\|Z+\Gamma Z\|_2^2\geq0.
\]
In particular $\ker T$ consists exactly of those elements of $L$
with $\Gamma Z=-Z$. Subtracting the orthogonal projection of $Y$
onto this kernel preserves (RT-OP-gap-equation). We can therefore
assume $Y\perp\ker T$. Partial transpose preserves the trace, so
(RT-OP-gap-equation) gives $\operatorname{Tr}Y=M^2$.

The projection $\Pi I$ is the low-isotypic orthogonal projector
$P_{<h}$ on $V\otimes V$. Its squared Hilbert--Schmidt norm is its
rank, namely $B$. Since $Y\in L$, Cauchy--Schwarz implies
\[
 M^4=|\langle Y,\Pi I\rangle|^2
       \leq B\|Y\|_2^2.
\]
It follows in particular that $B>0$. Projecting
(RT-OP-gap-equation) onto $L$ gives $TY=2\Pi I$ and hence
\[
 \langle Y,TY\rangle=2M^2.
\]
Thus $T$ has a positive eigenvalue. Since $Y\perp\ker T$,
\[
 \delta\frac{M^4}{B}
 \leq\delta\|Y\|_2^2
 \leq\langle Y,TY\rangle=2M^2,
\]
which proves (RT-OP-gap). The bounds
$M\geq q^n/(n+1)^{A_q}$ and
$B\leq(n+1)^{B_q}2^{nh}$ from
Proposition~\ref{pg:op-complete} and
Corollary~\ref{pg:op-rectangle} give the final inequality.
\end{proof}

For precision, $Q=(I+\Gamma)/2$ is the orthogonal projector onto
$E_+$, and $T=2\Pi Q\Pi|_L$. The positive eigenvalues of $T$ are
twice the squares of the nonzero singular values of $Q|_L$.
The quantity needed above is therefore the \emph{smallest positive}
one, not the largest overlap norm. The recoupling theorem of
Christandl, \c{S}ahino\u{g}lu, and Walter (2018, pinned author
preprint Section~2, Definition~1 and Eq.~(7), p.~4;
Lemma~2, p.~5; Section~3, Theorem~3, p.~6, proof pp.~6--8)
identifies projector-product operator norms with recoupling norms.
Its asymptotic statement gives an inverse-polynomial
Hilbert--Schmidt norm along some sequence for feasible quantum
marginals, and exponential decay along every sequence for infeasible
limiting marginals. It does not give the smallest-positive-singular-
value bound required by (RT-OP-gap). In addition, its existence
quantifier does not cover every prescribed rectangular type.
We therefore do not infer a high-support intersection, a square
price, or a shortening from that theorem.

\subsection{Why separate rank budgets did not close dimension four}

\begin{proposition}[RT-A3-D4-MOMENT-PROBE]
\label{rt:three-term-moment-probe}
The first-moment, second-moment, single-cut rank, and Lipschitz
conditions used in SH-58, without its mixed-cut inequalities,
admit an integer profile. Thus those four families alone do not
prove the dimension-four exclusion.
\end{proposition}
\begin{proof}
Index the eight complementary cuts by the binary masks $0,\ldots,7$.
Consider the following three rows:
\[
\begin{array}{c|rrrrrrrr|r|r}
 &0&1&2&3&4&5&6&7&\sum\kappa&\sum\kappa^2\\\hline
 1&12&24&12&24&12&12&12&24&132&2448\\
 2&12&12&24&12&24&12&24&12&132&2448\\
 3&24&12&12&9&12&24&0&12&105&1809
\end{array}
\]
All entries lie in $[0,24]$, every column includes a $24$, and
the last two columns are at most $141$ and $2496$, respectively.
For the folded distance
$h(s,t)=\min\{|s\triangle t|,4-|s\triangle t|\}$,
each row changes by at most $12$ across a distance-one edge;
distance-two differences are automatically at most $24$.
These assertions are also checked as exact integer inequalities
in \path{programs/rt_three_term_moment_profile_probe.py}, with
input and output \path{inputs/rt_three_term_moment_D4_probe.json}
and \path{outputs/rt_three_term_moment_D4_probe.json}.
The numerical search only proposes the integer rows; the retained
checker verifies all $1862$ integer constraints exactly.

For clarity, this profile fails a mixed-cut inequality already
used in SH-58. Take cuts $s=0$, $t=6$, at folded distance two,
and the SH-39 parameter $z=2048\leq4096=q$.
Then $p=4096$, and the three cheaper weighted rank bounds sum to
\[
 \sum_{j=1}^3\min\{2^{\kappa_{j,0}},z^2 2^{\kappa_{j,6}}\}
 =4096+4096+4194304=4202496
 <8388608=zp.
\]
It therefore cannot be the rank profile of a shortening identity.
This is a counterexample only to the sufficiency of the four
listed constraint families, and is neither a network layout nor
a tensor construction.
\end{proof}

\subsection{A sharp graph-certificate cost for the planar square}

\begin{proposition}[RT-SQUARE-GRAPH-SEED]
\label{rt:square-graph-seed}
Use the graph-certified pointwise averages of EX48 for
$P_2=\{\pm e_1/2,\pm e_2/2\}$. Their expected normalized
graph cost is at least $\sqrt2$. This lower bound is sharp for
the terminal-cut linear relaxation, even when arbitrary auxiliary
vertices, edge loads, and finite pointwise seed averages are allowed.
It is not an equality assertion for the actual spectral quantity
$\Gamma_2$.
\end{proposition}
\begin{proof}
Project each seed graph onto the target plane; projection preserves
the terminal-cut constraints and cannot increase edge lengths.
Normalize its edge loads by its positive terminal cut rate, so
every terminal-separating cut has capacity at least one. For
$u_\theta=(\cos\theta,\sin\theta)$, write
\[
 w_\theta(a)=\max_{p\in P_2}\langle u_\theta,a(p)\rangle
             -\min_{p\in P_2}\langle u_\theta,a(p)\rangle.
\]
Every threshold between the two extremes gives a terminal-separating
cut. Integrating its capacity over the threshold therefore gives
\[
 \sum_e\frac{b_e}{\rho}
       |\langle u_\theta,\Delta_e\rangle|
 \geq w_\theta(a).
\]
Integration over $0\leq\theta\leq\pi$ and
$\int_0^\pi|\langle u_\theta,v\rangle|\,d\theta=2\|v\|_2$
yield $2C(a)\geq\int_0^\pi w_\theta(a)\,d\theta$.
The width is a convex function of the labelled terminal positions:
the maximum is convex and the negative minimum is convex.
Thus a pointwise average with $\mathbb E a(p)=p$ satisfies
\[
 2\mathbb E C(a)\geq
 \int_0^\pi w_\theta(P_2)\,d\theta
 =\int_0^\pi\max\{|\cos\theta|,|\sin\theta|\}\,d\theta
 =2\sqrt2.
\]
For sharpness of the cut relaxation, place a cycle through the
four square vertices and assign load $\rho/2$ to each edge.
Every nontrivial terminal cut crosses at least two cycle edges.
Its normalized cost is half the square perimeter, namely
$\sqrt2$. No implication from cut feasibility to a tensor
conversion is needed for this sharpness statement.
\end{proof}

Consequently this graph-certificate class cannot establish a strict
improvement $\Gamma_2<\sqrt2$ merely by changing the planar graph
or averaging its terminal positions. EX48's cut-rank argument makes
the constraint necessary for graph-to-GHZ rate certificates over
every field. An improvement for actual spectral normalization would
require an argument outside this certificate class; the proposition
does not exclude such an improvement.
